\section{Discussion and Insights}

\textbf{What worked.} The pipeline is auditable. Hidden labels never reach the controller, every stop is
audited, and all trajectories are reported, including failures. No leakage was found by the automated
checks. The stop rule vetoed every stop at the one site where stopping would have failed. The TabPFN
hybrids were competitive, giving the best profile error and low-cell recall. The GP remained the most
conservative low-cell detector and the source of the uncertainty signal.

\textbf{What the results reveal.} Four lessons emerge.
\begin{enumerate}
\item \emph{Site predictability, more than the algorithm, decided whether CPTs could be saved.} Savings
  appeared only where five CPTs already describe the ground, and vanished at New Lock, where even 64 CPTs
  leave log-MAE near 0.9.
\item \emph{Location choice added nothing measurable here.} The benefit came from staged investigation
  combined with a stopping rule.
\item \emph{Model choice matters less than the decision it serves.} The five predictors lie within about
  3\% of each other on profile error. Their low-cell behaviour, however, differs sharply: GP is
  conservative, the graph hybrid sensitive and CNP blind. This difference changes what a stop rule can
  safely conclude.
\item \emph{Time, not CPT price, dominates the economics.} Per-batch decision delay can turn a neutral
  outcome strongly negative. A deployed system must return decisions within the working day, and should
  recommend the full plan early when a site is clearly unpredictable.
\end{enumerate}

\textbf{Threats to validity.}
\begin{itemize}
\item \emph{Optimistic development data.} The four development sites served as checkpoint-selection
  (validation) sites during source training. This likely explains why the learned models beat the GP by
  12\% on development sites but by only 1\% on evaluation sites. It also means the model ranking (0.578 vs
  0.597) is within noise.
\item \emph{Prior exposure.} The evaluation sites had been inspected in an earlier project stage, so this
  is a site-held-out re-evaluation, not a blind test.
\item \emph{Small samples.} With four evaluation sites, 5--16 audit positions per site and seeds that share
  the same ground, all conclusions are descriptive. The 7 observed early stops give a 95\% upper bound of
  roughly 40\% on the false-stop rate.
\item \emph{Confounding at New Lock.} New Lock is the only non-New-Zealand site, the largest site, and has
  no close source analogue. Domain shift, size and heterogeneity cannot be separated there.
\item \emph{Economic assumptions.} The economics rest on assumed unit values and on $\alpha=1$, which places
  investigation on the critical path.
\item \emph{Representation and labels.} The 0.25\,m grid misses thin layers, and the 1\,MPa threshold is a
  research indicator.
\end{itemize}

\textbf{Future improvements.} (i)~learn quality-predictive stop signals, rather than fixed quantile thresholds, so that the clear opportunities at Christchurch Stadium and Hollywood are taken; (ii)~add a site-predictability check after the initial CPTs that recommends the full plan at once for New Lock-like sites; (iii)~use larger batches or same-day decisions to cut latency; (iv)~use an event-weighted loss and a 0.1\,m grid for low-cell detection; (v)~evaluate on genuinely new sites with real contract and programme data.
A post-evaluation exploration (not used for any claim above) tried variance-reduction acquisition at New
Lock. It lowered the GP's budget-averaged profile error by about 11\% relative to maximin but worsened
thickness error, and none of 432 explored states met the targets.
